ここでは高被引用テーマの変化を、高被引用でない論文のそれと比較する。
各累積世代、各グループごとの最頻出被引用テーマと引用テーマを表\ref{tab:topmesh}に示す。これらのうちG1の被引用テーマ（高被引用テーマ）にのみにバラエティが確認される。
次に、各累積世代の被引用テーマと引用テーマの多様性を見るために各グループのエントロピーの変化を観察した。被引用テーマ、引用テーマ、被引用テーマと引用テーマ対の出現数のエントロピーを、それぞれ図\ref{fig:arr2}に示す。計数の方法は、被引用テーマと引用テーマについては、それぞれ論文ごとの分数カウントを採用する。被引用テーマ対引用テーマのペアについては、ひとつの参照ごとに被引用テーマと引用テーマの2部グラフを構成し、エッジの重みの合計が1となる分数カウントを採用する。図\ref{fig:arr2}からは高被引用論文のみがエントロピーが低く、その他のグループ間ではそれほど差がないことがわかる。なお、1946-1950年の累積世代に関してはMeSHディスクリプターの付与数が十分でなかったため除外している。
また、研究者の引用行動は必ずしも当該論文と同じテーマを引用するものではない。この転換の程度を見るために、グループごとに被引用テーマと引用テーマのコサイン距離を観察した。この結果を図\ref{fig:meshpivot}に示す。ここでも高被引用論文が累積世代の中期から後期にかけ他とは異なる様相を示した。また、高被引用論文は1946-2010年の累積世代以降、転換の程度は減少傾向にあり、他のグループとの差はなくなりつつあるように推察される。
